\subsection*{MT-M32. The corrected primary Euler quotient has a Riesz basis
and a local Poisson trace}

This result proves the local trace after the closure correction
of MT-M31.  It does not identify the global geometric kernel
with this local action.

\paragraph{Statement and conventions.}
Retain $H_s,w_s,U_h$ from MT-M31.  Fix $h>0$, positive
integers $d_1,\ldots,d_b$, and unit complex numbers
$a_1,\ldots,a_b$.  Define
\[
 Q(w)=\prod_{v=1}^{b}(1-a_vw^{d_v}),\qquad
 P=Q(U_h),\qquad
 Z_s=H_s/\overline{PH_s}^{\,H_s}.
 \tag{M32.1}
\]
Let $\lambda_1,\ldots,\lambda_L$ be the distinct roots
of $Q$, all on the unit circle, with multiplicities
$r_1,\ldots,r_L$.  Choose $0\leq\theta_\ell<2\pi$
such that $\lambda_\ell=e^{i\theta_\ell}$, and put
\[
 q_s=\max\{0,\lceil(s-1)/2\rceil\},\qquad
 r'_\ell=\min(r_\ell,q_s),\qquad
 t_{\ell,k}=\frac{\theta_\ell+2\pi k}{h}.
 \tag{M32.2}
\]
All roots and multiplicities here come from a finite Euler
polynomial; no Hecke GRH or Schanuel premise occurs.

The dual of $Z_s$, using the bilinear integral pairing of
MT-M31, has the Riesz basis
\[
 h^{-1/2}x^j e^{it_{\ell,k}x},
 \quad 1\leq\ell\leq L,\quad k\in\mathbf Z,
                  \quad0\leq j<r'_\ell.
 \tag{M32.3}
\]
When every $r'_\ell=0$, the quotient is zero and the basis
is empty.  The corresponding coefficient map on $Z_s$
is a bounded isomorphism onto the $\ell^2$ space of the
indices in (M32.3); its inverse supplies a biorthogonal
Riesz basis of the quotient.

For $f\in C_c^\infty(\mathbf R)$, or more generally for
an exponential-Schwartz test, let
\[
 C_f u(x)=\int_{\mathbf R}f(y)u(x-y)\,dy,
 \qquad M_f(z)=\int_{\mathbf R}f(y)e^{zy}\,dy.
 \tag{M32.4}
\]
The induced operator on $Z_s$ is trace class and satisfies
\[
 \operatorname{Tr}_{Z_s} C_f
      =\sum_{\ell=1}^{L}\sum_{k\in\mathbf Z}
                           r'_\ell M_f(it_{\ell,k})
      =h\sum_{n\in\mathbf Z}f(nh)
                          \sum_{\ell=1}^{L}r'_\ell\lambda_\ell^n.
 \tag{M32.5}
\]
Both expressions converge absolutely.  The map from tests
to trace-class operators is continuous in trace norm, with
uniform tails on bounded sets of tests at a fixed compact
support, and likewise for bounded sets in exponential
Schwartz space.

If $s>2\max_\ell r_\ell-1$, every local jet survives,
and the formula simplifies to
\[
 \operatorname{Tr}_{Z_s}C_f
       =\sum_{v=1}^{b}d_vh
                     \sum_{n\in\mathbf Z}a_v^{-n}f(nd_vh).
 \tag{M32.6}
\]
The strict condition is satisfied in particular by
$s\geq2b$.  For a primary number-field Euler product,
$h=\log p$, $d_v=f_v$ is the residue degree, and
$b\leq[F:\mathbf Q]$; hence the source's sufficient
degree $s\geq2[F:\mathbf Q]$ retains all multiplicities.

\paragraph{A finite matrix-weight model.}
If no index survives the cutoff, MT-M31 gives $Z_s=0$
and all the assertions are immediate.  Assume henceforth
that at least one index survives.
Let $\mathcal N_s\subset H_{-s}$ be the annihilator of
$PH_s$.  By the fully proved formula in MT-M31,
it is the closed span of the functions in (M32.3).
We construct a bounded isomorphism
\[
 \mathcal V_s:=\bigoplus_{\ell=1}^{L}
                     \bigoplus_{j=0}^{r'_\ell-1}L^2(0,h)
                  \longrightarrow\mathcal N_s.
 \tag{M32.7}
\]
For a coefficient vector $c=(c_{\ell,j})$, extend each
$c_{\ell,j}$ periodically with period $h$ and set
\[
 (Ac)(x)=\sum_{\ell,j}
       x^j e^{i\theta_\ell x/h}c_{\ell,j}(x\bmod h).
 \tag{M32.8}
\]
Endpoint representatives have measure zero and make no
difference to this definition.  Write $x=u+nh$ with
$0\leq u<h$, $n\in\mathbf Z$, and define
\[
 v_{\ell,j}(u,n)=(u+nh)^j
                         e^{i\theta_\ell(u/h+n)}.
 \tag{M32.9}
\]
The finite Hermitian matrix
\[
 G_{(\ell,j),(m,k)}(u)
   =\sum_{n\in\mathbf Z}
       \frac{\overline{v_{\ell,j}(u,n)}v_{m,k}(u,n)}
                  {(1+(u+nh)^2)^{s/2}}
 \tag{M32.10}
\]
converges absolutely and uniformly on $0\leq u\leq h$.
Indeed $j,k<q_s$ imply $j+k-s<-1$, so the tails
are bounded by a fixed summable power of $1+|n|$.
Thus $G(u)$ is continuous.  Tonelli, applied first to
finite coefficient sums and then to their squared absolute
values, gives the exact norm identity
\[
           \|Ac\|_{H_{-s}}^2
                =\int_0^h c(u)^*G(u)c(u)\,du.
 \tag{M32.11}
\]

The matrix $G(u)$ is positive definite for every fixed $u$.
If a vector $z$ has $z^*G(u)z=0$, then
\[
       \sum_{\ell,j}z_{\ell,j}
          e^{i\theta_\ell u/h}(u+nh)^j\lambda_\ell^n=0
                       \quad\text{for every }n\in\mathbf Z.
 \tag{M32.12}
\]
Exponential-polynomial sequences with distinct nonzero
$\lambda_\ell$ are linearly independent.
Here is a direct verification.  For one index $\ell$, apply
the finite-difference operator
$\prod_{m\ne\ell}(S-\lambda_m)^{r'_m}$, where
$(Sa)_n=a_{n+1}$, to (M32.12).  Every other component
vanishes.  If the $\ell$-component had a nonzero polynomial
of degree $j$, the resulting component would have the
same degree and a nonzero leading coefficient, multiplied
by $\prod_{m\ne\ell}(\lambda_\ell-\lambda_m)^{r'_m}$.
Applying $(S-\lambda_\ell)^j$ would then give a nonzero
constant multiple of $\lambda_\ell^n$, a contradiction.
Thus all the polynomials vanish.  The powers $(u+nh)^j$
are linearly independent polynomials in $n$, so $z=0$.

Continuity on the compact interval now gives constants
$0<c\leq C<\infty$, depending on this fixed local product
and $s$, with
\[
                  cI\leq G(u)\leq CI\quad(0\leq u\leq h).
 \tag{M32.13}
\]
Consequently (M32.8) is bounded and bounded below, and
its image is closed.  On each summand of $\mathcal V_s$,
the Fourier orthonormal basis
$h^{-1/2}e^{2\pi iku/h}$ maps to the corresponding
function in (M32.3).  Its image therefore contains a dense
subspace of $\mathcal N_s$ by MT-M31.  Conversely, finite
Fourier vectors are dense in $\mathcal V_s$ and their
images belong to the closed space $\mathcal N_s$, so
$A\mathcal V_s\subset\mathcal N_s$.  Its closed
image is exactly $\mathcal N_s$.  This proves (M32.7)
and the asserted Riesz basis, in the actual weighted norm.
No global Riesz assertion for a union of different rational
prime lattices is being made.

The bilinear integral pairing identifies $\mathcal N_s$
isometrically with the continuous dual of $Z_s$.
Taking the dual of the bounded isomorphism from the
standard $\ell^2$ basis onto (M32.3), and using reflexivity
of the Hilbert quotient, proves that the coefficient map
\[
 [u]\longmapsto
 \left(h^{-1/2}\int_{\mathbf R}u(x)x^je^{it_{\ell,k}x}\,dx
                                      \right)_{\ell,k,j}
 \tag{M32.14}
\]
is a bounded bijection with bounded inverse onto $\ell^2$.
This also proves the existence and uniform boundedness
of the biorthogonal quotient vectors.  It is stronger
than merely knowing that the individual coefficient
functionals are continuous.

\paragraph{The exact scalar quotient norm for one untwisted factor.}
For $P=1-U_h$ and $s>1$, the preceding model has a
particularly concrete primal form.  Put
\[
 W(u)=\sum_{n\in\mathbf Z}(1+(u+nh)^2)^{-s/2},
 \qquad
 (Rf)(u)=\sum_{n\in\mathbf Z}f(u+nh),\quad 0\leq u<h.
 \tag{M32.14a}
\]
The function $W$ is continuous and bounded above and
below by positive constants.  For almost every $u$,
Cauchy--Schwarz proves absolute convergence of $Rf(u)$
and
\[
 |Rf(u)|^2\leq W(u)
         \sum_{n\in\mathbf Z}|f(u+nh)|^2w_s(u+nh).
\]
Thus $R:H_s\to L^2((0,h),W(u)^{-1}du)$ is a contraction.
It is onto: a prescribed $F$ in the target has lift
\[
 f(u+nh)=\frac{F(u)}{W(u)w_s(u+nh)},
 \qquad
 \|f\|_{H_s}^2=\int_0^h\frac{|F(u)|^2}{W(u)}du.
 \tag{M32.14b}
\]
The coefficient $J_{2\pi k/h,0}(f)$ is the corresponding
Fourier coefficient of $Rf$ (with the positive-exponent
convention).  Completeness of Fourier series and MT-M31
give $\ker R=\overline{(1-U_h)H_s}$.  Therefore
\[
 H_s/\overline{(1-U_h)H_s}
        \simeq L^2((0,h),W(u)^{-1}du)
 \tag{M32.14c}
\]
is an isometry with its displayed quotient norm.
This is an explicit circle quotient even though the
actual difference-operator range is not closed.

\paragraph{The actual convolution and its trace norm.}
The polynomial-weight translation estimate of MT-M31
shows that $C_f$ is bounded on $H_s$ and $H_{-s}$.
It commutes with $P$ and preserves the closure of its range,
so it acts on $Z_s$.  On the dual functions one computes
directly
\[
 C_f^{\mathrm t}(x^je^{itx})
  =\int f(y)(x+y)^je^{it(x+y)}dy
  =\sum_{a=0}^{j}\binom ja M_f^{(a)}(it)
                                  x^{j-a}e^{itx}.
 \tag{M32.15}
\]
The same identity, paired with $u$, describes the quotient
operator in the coordinate map (M32.14).  It is the
direct sum over $(\ell,k)$ of triangular matrices of
size $r'_\ell$, with diagonal entry $M_f(it_{\ell,k})$.
Their sizes are bounded by the degree of $Q$, which is
fixed in this theorem.

For each integer $N\geq0$ and each derivative order
$0\leq a<\max r'_\ell$, integration by parts gives
\[
 |M_f^{(a)}(it)|\leq C_{N,a}(f)(1+|t|)^{-N}.
 \tag{M32.16}
\]
The constants are bounded by finitely many test seminorms:
on a fixed compact support they are controlled by finitely
many derivatives of $f$, and on exponential-Schwartz
space by finitely many of its defining seminorms.
For a finite $r\times r$ matrix its trace norm is at most
$r$ times its operator norm.  Hence the sum of the trace
norms of these blocks is finite, by (M32.16) with $N>1$.
Conjugating by the bounded isomorphism (M32.14) multiplies
the trace-norm estimate by at most the product of the two
isomorphism norms.  This proves trace class and its stated
continuity.

The block action was obtained on every coordinate of an
actual bounded convolution operator.  Thus the trace-class
sum equals that actual operator, first on finite coordinate
vectors and then by density.  Finite partial sums have
trace $\sum_{\ell,k}r'_\ell M_f(it_{\ell,k})$ over the
selected blocks.  Trace-norm convergence proves the first
identity in (M32.5).  The estimates also give uniform
tails on the stated bounded sets.  In particular no
unjustified passage from a Markushevich system to a
trace expansion is used.

\paragraph{Poisson summation with the normalization fixed.}
For a single $\theta$, periodize
$f(x)e^{i\theta x/h}$ with period $h$.
It is a smooth periodic function whose Fourier coefficients
are
\[
        h^{-1}M_f\bigl(i(\theta-2\pi k)/h\bigr).
\]
Its absolutely convergent Fourier series at zero gives
\[
 \sum_{k\in\mathbf Z}M_f\bigl(i(\theta+2\pi k)/h\bigr)
       =h\sum_{n\in\mathbf Z}f(nh)e^{i\theta n}.
 \tag{M32.17}
\]
The argument applies to exponential-Schwartz tests too,
since their periodization and all derivatives converge
uniformly.  Applying (M32.17) to each of the finitely
many roots proves the second equality in (M32.5).

If $r'_\ell=r_\ell$, count the roots with their
multiplicities in the original factorization of $Q$.
For a factor $1-a_vw^{d_v}$, its roots satisfy
$\lambda^{d_v}=a_v^{-1}$ and their power sum is
\[
 \sum_{\lambda^{d_v}=a_v^{-1}}\lambda^n
   =\begin{cases}
       d_va_v^{-n/d_v},&d_v\mid n,\\
       0,&d_v\nmid n.
     \end{cases}
 \tag{M32.18}
\]
This holds for positive, zero and negative integers $n$,
by the finite geometric-series identity on roots of unity.
Substituting in (M32.5) gives (M32.6).

\paragraph{Adelic normalization and geometric meaning.}
Here is an explicit conversion to a primary number-field
Euler product, with no hidden sign convention.
Choose a norm splitting
$C_F\simeq C_F^1\times\mathbf R$ and Haar measure
$dk\,dx$ with $\int_{C_F^1}dk=1$.
Fix a compact character $\beta$, extended to be trivial
on the chosen norm section, and write vectors in its
sector as $\Phi(k,x)=\beta(k)^{-1}u(x)$.
For a finite place $v\mid p$, write the uniformizer idele
as
\[
             \pi_v=(k_v,-f_v\log p).
\]
The literal operator in (M31.1) becomes
$u(x)-\beta(k_v)u(x+f_v\log p)$.
After the isometric reflection $x\mapsto-x$ this is
the factor in (M32.1) with
\[
              h=\log p,\quad d_v=f_v,
                         \quad a_v=\beta(k_v).
 \tag{M32.19}
\]
For a Bruhat test $\varphi\in\mathcal D(C_F)$, put
\[
 b_\beta(y)=\int_{C_F^1}\varphi(k,y)\beta(k)\,dk,
 \qquad
 \varphi_\beta(k,y)=\beta(k)^{-1}b_\beta(y).
 \tag{M32.20}
\]
The actual convolution on this sector is $C_{b_\beta}$;
after reflection its test is $f(y)=b_\beta(-y)$.
Consequently, when all primary jets survive, (M32.6) reads
\[
     \operatorname{Tr} A_\beta(\varphi)
       =\sum_{v\mid p}\log q_v
                          \sum_{n\in\mathbf Z}\varphi_\beta(\pi_v^n).
 \tag{M32.21}
\]
Only finitely many $n$ contribute for a compact radial
support.  If $\beta$ is unramified at $v$,
$\varphi_\beta$ is constant on the local unit group.
For any multiplicative Haar measure $\mu_v$, the
corresponding local expression is equivalently
\[
 \frac{\log q_v}{\mu_v(\mathcal O_{F_v}^{\times})}
          \int_{F_v^{\times}}\varphi_\beta(g_v)\,d\mu_v(g_v).
 \tag{M32.22}
\]
This specifies exactly where the local-unit volume
normalization enters.  The formulas do not silently
replace an integral by a sum while retaining its volume factor.

The finite matrix $G(u)$ is a concrete Hilbert model for
the finite local Euler polynomial $Q$, with polynomial
jet degrees retaining its nilpotent extensions.  Fourier
series on the norm circle supply the infinite lattice of
Euler poles.  The degree bound used for the local trace
therefore has a direct periodic-function explanation.
It supplies the corrected local quotient and trace intended
by Masullo Theorem 11.2.2.1, original pp.~490--491,
after MT-M31's closure correction and the strict jet cutoff.

The constants in (M32.13) are for a fixed prime product,
character and degree.  No uniform comparison as primes,
characters or global geometric truncations vary is asserted.
In particular this result does not show that the
localized kernel in Masullo Theorem 11.3.2.1,
pp.~495--496, acts on this corrected quotient in the
claimed way, nor that its nuclear limits commute with
the global shell degeneration.  That precise comparison
remains the subsequent geometric-trace gate.
